\chapter{\MakeUppercase{Metodologia}}\label{sec:metodologia}

\section{Delineamento e coleção experimental}\label{sec:met-delineamento}
\label{sec:met-instancias}

A investigação foi organizada em torno da Análise do Espaço de Instâncias
(ISA), aplicada a uma coleção de instâncias do WSP e ao desempenho de seis
abordagens (MIP, LBBD, ILS, IBS, CUTH e RS). Cada instância foi a unidade de
análise; as execuções com diferentes sementes forneceram as repetições,
agregadas por instância. O estudo começou pela exploração das 55 instâncias
originais e pela construção de um mapa MATILDA de referência, mantido fixo.
Em seguida, aplicou-se a geração orientada de instâncias: definiram-se
pontos-alvo em regiões pouco ocupadas desse mapa, e um algoritmo genético,
inspirado em \citeonline{SmithMilesBowly2015}, buscou configurações do gerador
de \citeonline{DelazeriRitt2026} cujas instâncias se projetassem próximas
desses alvos. As 19 instâncias aceitas ampliaram a coleção para 74,
utilizada na avaliação das seis abordagens e no reajuste dos quatro modelos ISA,
com e sem medidas topológicas adicionais. Como etapa complementar,
registraram-se as trajetórias de busca do IBS, conforme a
Seção~\ref{sec:met-trajetorias}.
A Figura~\ref{fig:met-fluxo} sintetiza essa organização.

\begin{figure}[htbp]
\centering
\caption{Etapas de caracterização, geração e análise experimental.}
\label{fig:met-fluxo}
\begin{singlespace}
\begin{tikzpicture}[
  etapa/.style={draw, align=center, text width=8.2cm, minimum height=0.85cm,
    inner sep=5pt, font=\small},
  fluxo/.style={-{Stealth[length=2mm]}, thick}]
\node[etapa] (base) at (0,0) {Coleção inicial: 55 instâncias};
\node[etapa,text width=3.6cm] (historicas) at (6.8,-1.5) {STNs históricas:\\16 instâncias clássicas;\\160 execuções};
\node[etapa] (atributos) at (0,-1.5) {Caracterização original e topológica\\e análise exploratória};
\node[etapa] (mapa) at (0,-3) {Mapa MATILDA de referência\\povoamento e cobertura};
\node[etapa] (geracao) at (0,-4.5) {Povoamento orientado, validação\\e seleção de 19 instâncias adicionais};
\node[etapa] (colecao) at (0,-6) {Coleção ampliada: 74 instâncias};
\node[etapa,text width=5.4cm] (benchmark) at (-3.2,-7.5) {Comparação dos seis métodos\\nas 74 instâncias};
\node[etapa,text width=5.4cm] (trajetorias) at (3.2,-7.5) {Registro de trajetórias do IBS\\nas 19 instâncias adicionais};
\node[etapa,text width=5.4cm] (isa) at (-3.2,-9.15) {Modelos finais MATILDA:\\controle/topologia, 55/74 instâncias};
\node[etapa,text width=5.4cm] (redes) at (3.2,-9.15) {Construção das STNs\\e cálculo dos descritores};
\node[etapa] (analise) at (0,-11.3) {Análise comparativa dos desempenhos,\\das representações e das trajetórias};
\draw[fluxo] (base) -- (atributos);
\draw[fluxo] (base.east) -- ++(0.25,0) |- (historicas.west);
\draw[fluxo] (atributos) -- (mapa);
\draw[fluxo] (mapa) -- (geracao);
\draw[fluxo] (geracao) -- (colecao);
\draw[fluxo] (colecao) -- (benchmark);
\draw[fluxo] (colecao) -- (trajetorias);
\draw[fluxo] (benchmark) -- (isa);
\draw[fluxo] (trajetorias) -- (redes);
\draw[fluxo] (isa.south) -- (analise.north);
\draw[fluxo] (redes.south) -- (analise.north);
\draw[fluxo] (historicas.east) -- (9.05,-1.5) -- (9.05,-11.3) -- (analise.east);
\end{tikzpicture}
\par\smallskip
{\fonte{elaboração própria.}}
\end{singlespace}
\end{figure}

As entradas vieram da coleção de instâncias associada aos trabalhos de
\citeauthoronline{DelazeriRitt2024} (\citeyear{DelazeriRitt2024,DelazeriRitt2026}). A consolidação em ordem
determinística reteve a primeira ocorrência de cada identificador:
16 casos clássicos LA0--LA7 e LB0--LB7 e 39 do gerador. A geração orientada
acrescentou 19 casos, conforme a Seção~\ref{sec:met-geracao}.
A coleção constituiu um conjunto dirigido de casos de teste do modelo computacional.

A classe Small reuniu casos com até 400 vértices; Medium, Large e Huge
corresponderam a 900, 1.600 e 6.400 vértices. A
Tabela~\ref{tab:met-colecao} apresenta composição, repetições e orçamentos
para a comparação dos seis métodos.

\begin{table}[htbp]
\centering
\caption{Composição da coleção e protocolo da comparação dos seis métodos.}
\label{tab:met-colecao}
\normalsize
\begin{tabular}{lrrrrr}
\hline\hline
Classe & Originais & Adicionais & Total & Repetições & Limite (s) \\
\hline
Small  & 19 & 1 & 20 & 10 & 600 \\
Medium & 30 & 3 & 33 & 10 & 1.350 \\
Large  & 3 & 6 & 9 & 10 & 2.400 \\
Huge   & 3 & 9 & 12 & 3 & 9.600 \\
\textbf{Total} & \textbf{55} & \textbf{19} & \textbf{74} & -- & -- \\
\hline\hline
\end{tabular}
\par\smallskip
{\singlespacing\fonte{elaboração própria a partir do planejamento experimental.}}
\end{table}

Cada arquivo no formato \textit{JavaScript Object Notation} (JSON) descreveu
o grafo, os tempos dos arcos, as coordenadas, a ignição, o horizonte e os
vetores de liberações, capacidades e atrasos. A validação das entradas
considerou as dimensões do grafo, a unicidade dos identificadores e dos
arcos, a positividade e a finitude dos tempos e a compatibilidade dos vetores.
Seu comprimento definiu o número de pontos de liberação; a soma das
capacidades forneceu o total de recursos. Nas 74 instâncias, os valores do
vetor de atrasos são iguais dentro de cada instância e correspondem ao
parâmetro $\Delta$ do modelo, embora possam variar entre instâncias.

As entradas de ILS e LBBD foram convertidas para o formato utilizado nas
implementações de \citeonline{HarrisForbesTaimre2023}, preservando ignição, arcos, tempos
decimais, horizonte, atrasos e capacidades. Recursos simultâneos foram
agregados quando necessário. A equivalência foi verificada por comparação
semântica e avaliação das soluções; também se verificou a identidade dos
arquivos de entrada utilizados nos dois computadores do experimento.

\section{Caracterização original e topológica}\label{sec:met-topologia}

Extraíram-se 36 atributos originais, agrupados em tamanho e estrutura,
recursos e liberações, tempos dos arcos, coordenadas espaciais e propagação
sem intervenções. Incluíram-se $|V|$, $|A|$, sua razão, horizonte,
capacidade total, pontos de decisão e recursos por ponto; atraso e sua
razão pelo horizonte; e primeira e última liberações e janela temporal,
normalizadas pelo horizonte. A razão $|A|/|V|$ corresponde ao número médio
de arcos por vértice, não à densidade usual normalizada por $|V|(|V|-1)$.

A propagação livre foi calculada por Dijkstra. Seus descritores incluíram
maior chegada finita sem recursos, razão pelo horizonte, fração queimada antes dele e
quantis de 5\%, 50\% e 95\% normalizados. Também foram consideradas
assimetria dos tempos recíprocos, amplitudes espaciais, dispersão vertical
e estatísticas dos pesos. Nos instantes de liberação, o perímetro contou
vértices atingidos com vizinhos ainda não atingidos; o número de chegadas
entre liberações (janela frontal) contou os vértices atingidos entre
liberações sucessivas, ou até o horizonte na última.
Para ambos, extraíram-se média e máximo. Os atributos das 55 originais
foram recalculados na ampliação e conferidos com tolerância $10^{-9}$.
A Tabela~\ref{tab:met-atributos-originais} detalha os 36 atributos.

\begin{table}[p]
\centering
\caption{Definição dos 36 atributos originais extraídos por instância.}
\label{tab:met-atributos-originais}
\begin{singlespace}
\setlength{\tabcolsep}{4pt}
\renewcommand{\arraystretch}{1.12}
\begin{tabular}{>{\raggedright\arraybackslash}p{2.35cm}r>{\raggedright\arraybackslash}p{9.8cm}}
\hline\hline
\textbf{Grupo} & \textbf{Qtd.} & \textbf{Atributos e cálculo}\\
\hline
Dimensão e horizonte & 4 & $|V|$, $|A|$, $|A|/|V|$ e $H$.\\
Recursos e atraso & 5 & Total $k=\sum_i q_i$; número $T$ de liberações; $k/T$; atraso médio, igual a $\Delta$ nestas instâncias; e $\Delta/H$.\\
Liberações & 3 & Primeira liberação/$H$ ($\tau_1/H$), última liberação/$H$ ($\tau_T/H$) e janela de liberação/$H$ ($(\tau_T-\tau_1)/H$). Esta última é o intervalo entre a primeira e a última liberação, relativo ao horizonte.\\
Tempos dos arcos & 8 & Média, desvio padrão populacional, mínimo, máximo, quantis de 25\%, 50\% e 75\%, e coeficiente de variação (desvio dividido pela média) de $t_{uv}$.\\
Assimetria & 2 & Média e máximo de $|t_{uv}-t_{vu}|/[(t_{uv}+t_{vu})/2]$, contando uma vez cada par de arcos recíprocos; zero se não houver pares.\\
Coordenadas & 4 & Amplitudes das coordenadas $x$, $y$ e $z$ e desvio padrão populacional de $z$.\\
Chegadas sem recursos & 6 & Maior chegada finita sem recursos e sua razão por $H$; fração dos vértices com chegada menor que $H$; percentil 5, mediana e percentil 95 das chegadas finitas, divididos por $H$.\\
Perímetro nas liberações & 2 & Média e máximo do número de vértices queimados antes de $\tau_i$ com ao menos um sucessor ainda não queimado.\\
Chegadas entre liberações & 2 & Média e máximo de chegadas entre liberações: número de vértices atingidos em $[\tau_i,\tau_{i+1})$; na última liberação, em $[\tau_T,H)$. Esse descritor também é denominado janela frontal.\\
\hline
\textbf{Total} & \textbf{36} & \\
\hline\hline
\end{tabular}
\par\smallskip
{\small Nota: chegadas e perímetro são calculados sem recursos, por Dijkstra.
A janela de liberação é um intervalo temporal; chegadas entre liberações é uma contagem de vértices.
O perímetro é um descritor de fronteira queimada; difere da
janela de candidatos ainda não queimados usada pelo IBS. A fração queimada
livre, constante na base, é retirada da entrada final da ISA.\par}
{\fonte{elaboração própria a partir das definições implementadas na extração dos atributos.}}
\end{singlespace}
\end{table}

A exploração inicial utilizou os desempenhos históricos das seis abordagens
nas 55 instâncias originais
(\citeauthor{DelazeriRitt2024}, \citeyear{DelazeriRitt2024,DelazeriRitt2026}),
agregados por instância e família. Quando um método dispunha de mais de uma
variante nesses dados, como o IBS, considerou-se em cada instância a de menor
mediana de \textit{gap}, de modo a representar o melhor resultado histórico
do método. Esses desempenhos serviram apenas a esta etapa exploratória; o
experimento final adotou uma configuração fixa por método, descrita na
Seção~\ref{sec:met-benchmark}. Após imputação pela mediana e padronização,
aplicaram-se PCA bidimensional e florestas aleatórias para explorar a relação
dos atributos com o melhor método e a menor mediana de \textit{gap}.

A topologia acrescentou 17 medidas nodais: graus de entrada e saída,
proporção de entrada e desequilíbrio entre graus; PageRank; hubs e
authorities do HITS; intermediação; harmônicas de entrada e saída;
e agrupamento direcionado. PageRank, HITS, intermediação e harmônicas
foram calculados com e sem ponderação. Cada medida foi resumida por
média, desvio padrão, mínimo, quantis de 25\%, 50\%, 75\% e 90\%,
e máximo, totalizando 136 candidatos. Os resumos permitiram relacionar
a topologia de cada instância ao desempenho dos algoritmos.

Os tempos dos arcos foram usados como distâncias nos caminhos mínimos;
nas medidas em que peso expressa força de conexão, PageRank e HITS,
utilizaram-se seus inversos, conforme a Seção~\ref{sec:topologia}.
O PageRank adotou amortecimento 0,85, valor convencional associado à
proposta de \citeonline{BrinPage1998}. A tolerância de $10^{-12}$ e o
limite de 1.000 iterações foram escolhas numéricas do estudo, para exigir
convergência precisa e permitir a conclusão do cálculo nos grafos maiores.
Esses valores não foram ajustados para maximizar associações com desempenho.
Os escores de PageRank foram normalizados para soma um.

O HITS foi obtido por decomposição esparsa em valores singulares, repetida
com sementes 11, 23, 37, 53 e 71 para examinar a dependência do estado inicial.
Os escores de cada execução foram normalizados para soma um e agregados
pela média. Amplitudes superiores a $10^{-6}$ entre sementes indicaram
instabilidade. Esse limiar absoluto é uma tolerância operacional para
separar variações de arredondamento de diferenças maiores nos escores
normalizados, e não um teste estatístico ou uma constante teórica do HITS.
Sua sensibilidade a outros valores não foi avaliada.

Intermediação e centralidades harmônicas foram calculadas exatamente,
sem amostragem de pares, com igraph. O uso de rotinas compiladas tornou
viável manter o cálculo exato nos grafos maiores e evitou que erro de
amostragem alterasse os atributos usados na projeção das candidatas.
As normalizações foram, respectivamente, $(|V|-1)(|V|-2)$ e $|V|-1$,
com contribuição nula de pares inalcançáveis, permitindo comparar tamanhos
distintos. PageRank e agrupamento direcionado de \citeonline{Fagiolo2007}
utilizaram NetworkX, mantendo as definições adotadas para essas medidas.
A escolha das bibliotecas foi computacional; não constituiu uma comparação
de métodos topológicos ou de pacotes.

A auditoria excluiu atributos constantes ou com variância numericamente
nula, pois não discriminam instâncias; duplicatas exatas, por redundância;
e atributos derivados de HITS instável, para não tratar variações numéricas
como diferenças do problema. Também excluiu degenerações no deslocamento
positivo do PRELIM: valores distintos em precisão de máquina podiam se
tornar idênticos após subtrair o mínimo e somar um, inviabilizando Box--Cox
e a padronização posterior.

Dos 136 candidatos topológicos, 34 foram classificados como constantes
ou degenerados, incluindo os oito resumos do agrupamento; 16 foram
excluídos por HITS instável, 11 por duplicação exata e dois por degeneração
numérica no PRELIM. Restaram 73, mantidos nas coleções de 55 e 74 instâncias.
Dos 36 originais, retirou-se a fração queimada livre, constante, formando
o controle com 35. Na combinação, removeu-se a média do grau de entrada,
idêntica a $|A|/|V|$, pois a soma dos graus de entrada é $|A|$.
Assim, a entrada combinada teve $35+73-1=107$ candidatos. A preparação
exigiu matrizes completas e finitas, com desempenho medido para todas
as combinações utilizadas.

\section{Geração orientada e validação das adições}\label{sec:met-geracao}

A geração orientada usou como referência o mapa MATILDA das 55 instâncias
originais, construído com os atributos originais e topológicos e com os
desempenhos históricos das seis abordagens. Esse mapa foi mantido fixo durante
toda a geração, de modo que alvos, projeções e medidas de cobertura se
referissem sempre à mesma representação. A evolução em direção aos alvos
seguiu a estratégia de \citeonline{SmithMilesBowly2015}.

O mapa foi ajustado a partir de 108 atributos candidatos (35 originais e 73
topológicos), dos quais se selecionaram dez. A média do grau de entrada,
redundante com $|A|/|V|$, só foi retirada dos 107 candidatos dos modelos
finais (Seção~\ref{sec:met-topologia}). Foram considerados bons os
desempenhos com \textit{gap} de até 10\% (limiar 0,10); os dados passaram por
Box--Cox e padronização; o PILOT projetou as instâncias em duas dimensões,
com cinco tentativas de otimização; e o PYTHIA ajustou os classificadores
com cinco partições estratificadas. O número de partições foi limitado pela
classe menos numerosa: nos rótulos históricos, apenas cinco instâncias não
atingiram o critério de bom desempenho para o IBS, de modo que cinco
partições estratificadas ainda permitem exemplos das duas classes em cada
uma. Os modelos finais, ajustados com os desempenhos do experimento
comparativo, têm configuração própria (Seção~\ref{sec:met-matilda}).

Os alvos foram escolhidos em três etapas, de modo que fossem viáveis,
ficassem em regiões pouco cobertas e se distribuíssem pelo mapa. Primeiro,
delimitou-se o domínio dos alvos como a interseção entre o fecho convexo dos
55 pontos do mapa e a região plausível estimada pelo CLOISTER
(Seção~\ref{sec:matilda}), evitando alvos fora da região observada ou em
combinações de atributos tidas como inviáveis. Nesse domínio, uma sequência
de Sobol bidimensional embaralhada (semente 20260612,
Quadro~\ref{qua:met-geracao}) forneceu 10.000 pontos admissíveis,
distribuídos de forma quase uniforme.

{\emergencystretch=2em\relax
Segundo, identificaram-se os pontos pouco cobertos. Para cada instância
original, calculou-se a distância no mapa até a instância mais próxima,
excluindo ela mesma; o quantil de 90\% dessas distâncias definiu o limiar de
subcobertura. Foram considerados pouco cobertos os pontos candidatos cuja
distância à instância mais próxima excedia esse limiar, isto é, mais isolados
do que 90\% das instâncias existentes.\par}

Terceiro, entre os pontos pouco cobertos, selecionaram-se os alvos pelo
critério maximin: escolhe-se o ponto mais distante das instâncias já
presentes, que passa a ser tratado como ocupado, e o processo se repete.
Isso espalha os alvos pelo mapa, evitando concentrá-los na mesma lacuna.

Cada uma das três sementes tratou dois lotes de seis alvos. O primeiro
foi comum; o segundo foi adaptativo, considerando as aceitações de cada
semente e nova amostra Sobol. Os alvos adaptativos foram analisados
segundo a semente e o lote de seleção. Cada indivíduo codificou
oito fatores do gerador e uma semente: tamanho, inclinação, vento, atraso,
quantidade de recursos e de pontos de liberação, primeira e última
liberações. Os domínios categóricos originais foram preservados.

A aptidão minimizou a distância euclidiana ao alvo. A projeção da candidata
reutilizou deslocamento positivo, Box--Cox, padronização, seleção e matriz
PILOT do mapa de referência, sem limitação de extremos. A projeção foi
conferida nos 55 pontos com tolerância $10^{-8}$. Durante a busca,
calcularam-se as famílias topológicas necessárias à seleção; os finalistas
receberam extração completa e recomputação da projeção. O gerador utilizou
Dijkstra com fila de prioridade, conservando a regra de caminhos mínimos.
O Quadro~\ref{qua:met-geracao} apresenta a configuração; os 36 alvos
consumiram 72.000 avaliações, repartidas entre evolução e exploração local.

\begin{quadro}[htbp]
\centering
\caption{Configuração da geração orientada de instâncias.}
\label{qua:met-geracao}
\begin{singlespace}
\normalsize
\begin{tabular}{|p{5.0cm}|p{8.0cm}|}
\hline
\textbf{Componente} & \textbf{Configuração adotada} \\
\hline
Sementes principais & 20260612, 104729 e 1299709 \\
\hline
Alvos por semente & 12, organizados em dois lotes de seis \\
\hline
População e evolução & 40 indivíduos; 40 gerações \\
\hline
Seleção e elitismo & Torneio de tamanho três; dois indivíduos preservados \\
\hline
Cruzamento e mutação & Uniforme, com probabilidade 0,8; mutação independente por gene, com probabilidade $1/9$ \\
\hline
Orçamento por alvo & 1.600 avaliações genéticas e 400 locais; total de 2.000 \\
\hline
Aptidão & Minimização da distância euclidiana ao alvo no mapa fixo \\
\hline
\end{tabular}
\par\smallskip
{\fonte{elaboração própria a partir da configuração da geração.}}
\end{singlespace}
\end{quadro}

Os valores do Quadro~\ref{qua:met-geracao} são configurações da adaptação
ao WSP, registradas no experimento; não constituem reprodução integral dos
parâmetros de \citeonline{SmithMilesBowly2015}. A inspiração desse trabalho
é o uso da distância ao alvo no espaço de instâncias como objetivo da
evolução. A codificação dos fatores do gerador, os operadores, a população
e os orçamentos foram definidos para esta aplicação. A divisão de 80\%
das avaliações para evolução e 20\% para exploração local segue uma
adaptação da estratégia híbrida discutida por \citeonline{Bowly2019},
com gerações síncronas e processos independentes por semente. Não houve
comparação sistemática para determinar parâmetros ótimos para o WSP.

Seja $\rho$ a distância do alvo ao ponto ocupado mais próximo no início do
lote. Uma candidata foi \textit{hit} se sua distância foi até $0{,}25\rho$;
\textit{improved} se superior a $0{,}25\rho$ e inferior a $\rho$; e
\textit{missed} se não melhorou ou não passou pela validação.
A aceitação exigiu consistência, projeção finita e pertencimento ao
domínio interno após recomputação exata, com tolerância $10^{-10}$.
A união das aceitações das três sementes foi deduplicada: arquivos de
conteúdo idêntico seriam contados uma única vez, mediante hash SHA-256 do
JSON em forma canônica, com chaves ordenadas e formatação padronizada.
As 19 aceitações, sete, seis e seis por semente, já eram distintas;
a deduplicação manteve 19 arquivos únicos.

A cobertura foi avaliada no mesmo mapa fixo e nos mesmos 10.000 pontos
Sobol. Mediram-se maior distância à coleção (\textit{fill distance}),
média, mediana, quantil de 90\% e fração acima do limiar de subcobertura.
Calculou-se também a proporção de células ocupadas em grades de
$8\times8$, $10\times10$ e $12\times12$ sobre o retângulo delimitador
do fecho convexo; os resultados apresentam a grade $10\times10$.
As sequências foram avaliadas separadamente sobre a base original.
A comparação entre a coleção original e a ampliada manteve fixos o domínio
de análise e a transformação para o espaço de instâncias, permitindo
quantificar o preenchimento das lacunas no mapa de referência.

\section{Experimento comparativo e ambiente}\label{sec:met-benchmark}

IBS, MIP, CUTH e RS utilizaram implementações em C++ associadas aos
trabalhos de \citeauthoronline{DelazeriRitt2024} (\citeyear{DelazeriRitt2024,DelazeriRitt2026}); MIP correspondeu ao modelo básico associado
a \citeonline{DelazeriRitt2026}. ILS e LBBD utilizaram a implementação
Python de \citeonline{HarrisForbesTaimre2023}, com controle
do tempo de execução e validação do objetivo. Os parâmetros estão no
Quadro~\ref{qua:met-algoritmos}.

\begin{quadro}[htbp]
\centering
\caption{Configurações dos métodos no experimento comparativo.}
\label{qua:met-algoritmos}
\begin{singlespace}
\normalsize
\begin{tabular}{|p{1.5cm}|p{11.5cm}|}
\hline
\textbf{Método} & \textbf{Configuração adotada} \\
\hline
IBS & $\beta=50$, $\eta=70$, $c=30$, $p=0{,}5$, $\hat{p}=1$ e $z_{\max}=3$; sem limite adicional de iterações \\
\hline
MIP & Modelo básico; Gurobi 10.0.3, uma \textit{thread}; limite global por classe; \texttt{Symmetry=-1}, \texttt{RINS=-1}, \texttt{MIPFocus=0}, \texttt{Heuristics=0.05}, \texttt{ImproveStartGap=0} e \texttt{ImproveStartTime} sem gatilho efetivo; pré-processamento habilitado \\
\hline
CUTH & Gurobi 10.0.3, uma \textit{thread}; tolerância absoluta de 0,01 na busca binária temporal; até 300 segundos por chamada ao modelo de corte, limitados pelo tempo global restante \\
\hline
RS & Semente da execução; limite global por classe \\
\hline
ILS & Até 50 reinícios na construção inicial; limites de cinco candidatos, 20 vizinhos, cinco modificações, 100 falhas e 50 passos sem melhoria; $p_1=0{,}075$ e $p_2=0{,}025$ \\
\hline
LBBD & Inicialização heurística e etapa exata sob um orçamento compartilhado; Gurobi 10.0.3, uma \textit{thread}; cortes por \textit{lazy constraints} \\
\hline
\end{tabular}
\par\smallskip
{\small Nota: Gurobi é um \textit{solver} comercial de otimização matemática. A MIP, a LBBD e o submodelo da CUTH utilizaram a versão 10.0.3. Os parâmetros listados para a MIP vieram da implementação empregada, sem ajuste neste estudo.\par}
{\fonte{elaboração própria a partir dos parâmetros do experimento comparativo.}}
\end{singlespace}
\end{quadro}

O IBS manteve os parâmetros de \citeonline{DelazeriRitt2024}, com
$\hat p=1$ conforme \citeonline{DelazeriRitt2026}, sem limite adicional
de iterações. A CUTH utilizou os valores padrão da implementação.
Os mecanismos dos parâmetros de IBS e ILS foram descritos nas
Seções~\ref{sec:ibs} e \ref{sec:ils}; seus valores no experimento estão
no Quadro~\ref{qua:met-algoritmos}. As sementes foram repassadas aos
procedimentos que admitiam esse controle.

Na ILS, o relógio monotônico controlou o prazo, admitindo para pontuação
somente candidatos completos produzidos dentro dele. Na LBBD,
inicialização e etapa exata compartilharam o orçamento. Sem incumbente
utilizável, retornou-se a alternativa viável sem recursos. Objetivo,
limites e estado do modelo foram preservados separadamente da pontuação
reavaliada no modelo nativo, conforme a Seção~\ref{sec:met-analise}.

Os limites da Tabela~\ref{tab:met-colecao} seguiram o orçamento de
1,5 segundo por célula de \citeonline{DelazeriRitt2026};
os casos clássicos também receberam 600 segundos.
Usaram-se sementes 1--10 em Small, Medium e Large e 1--3 em Huge,
totalizando 656 execuções por método e 3.936 na comparação.
As 190 trajetórias instrumentadas formaram uma fase separada.

O experimento foi dividido entre dois computadores com Intel Core i7-7700
3,60 GHz, quatro núcleos físicos, oito processadores lógicos e cerca de
16 GB de memória. Os ambientes registraram Linux x86\_64, kernel
\texttt{7.0.0-31-generic}, Gurobi/gurobipy 10.0.3, Python 3.11.16
na otimização e 3.12.14 na preparação. Cada máquina recebeu 328 trabalhos
por método e 95 trajetórias, sem duplicar pares instância--semente.
Processos foram fixados em núcleos físicos distintos, com até quatro
trabalhos sem \textit{solver} ou dois dependentes de Gurobi e bibliotecas numéricas
restritas a uma thread. Entradas e executáveis foram confrontados por hashes.

Registraram-se solução, objetivo, status, tempo de parede e pico de memória.
O tempo de busca e a margem externa para encerramento e gravação foram
registrados separadamente. Caracterização e geração utilizaram outro ambiente, com Python
3.12.3, quatro processadores lógicos e cerca de 8 GB de memória.
Os ajustes MATILDA usaram Python 3.12.3, dois trabalhadores por ajuste e
uma implementação do fluxo baseada em \texttt{pyInstanceSpace} 0.2.1.
As configurações efetivamente utilizadas foram registradas junto às
entradas e saídas de cada ajuste.
As bibliotecas utilizadas incluíram NumPy 1.26.4, pandas 2.2.3, SciPy
1.14.1, scikit-learn 1.5.2, PyGAD 3.7.0, NetworkX 3.4.2 e igraph 1.0.0.

\section{Reconstrução da ISA e avaliação das representações}\label{sec:met-matilda}

Construíram-se quatro modelos MATILDA para comparar as coleções e as
representações: o controle, com os atributos originais, e a combinação
desses atributos com medidas topológicas adicionais. Assim, controle não
significa ausência de estrutura do grafo, pois $|V|$, $|A|$ e sua razão
já pertencem aos atributos originais. As entradas de desempenho das 55
instâncias foram subconjuntos exatos das 74 linhas de \textit{gaps} auditados. Cada matriz teve seis colunas de
algoritmos e 35 atributos candidatos no controle ou 107 com topologia.
A contribuição preditiva foi examinada por diagnósticos internos de
classificação dos rótulos de desempenho. Utilizaram-se máquinas de vetores
de suporte (SVM, do inglês: \textit{support vector machines}) com núcleo
de função de base radial (RBF, do inglês: \textit{radial basis function}).

O desempenho de entrada foi a média do \textit{gap} relativo, em fração. O rótulo
bom admitiu objetivo médio até 10\% acima do melhor valor conhecido
(BKV, do inglês: \textit{best known value}). Esse critério foi mantido
nos quatro modelos e aplicado por algoritmo, permitindo vários métodos
bons na mesma instância. O PRELIM aplicou
deslocamento positivo, Box--Cox e padronização, sem redução de instâncias
ou limitação de extremos. As configurações estão no
Quadro~\ref{qua:met-matilda}.

\begin{quadro}[htbp]
\centering
\caption{Parâmetros do fluxo MATILDA nos quatro modelos.}
\label{qua:met-matilda}
\begin{singlespace}
\normalsize
\begin{tabular}{|p{2.2cm}|p{10.8cm}|}
\hline
\textbf{Etapa} & \textbf{Configuração adotada} \\
\hline
PRELIM & Deslocamento positivo, Box--Cox e padronização; sem limitação de valores extremos ou redução de instâncias \\
\hline
SIFTED & Correlação mínima 0,1 e significância 0,05; dez grupos, 100 inicializações e até 1.000 iterações do agrupamento \\
\hline
Seleção genética & População 50; até 100 gerações; parada após cinco gerações sem melhoria; torneio de tamanho quatro, dois pais e elitismo de dois; cruzamento \textit{scattered} 0,8; mutação aleatória 0,05; semente de base 0 \\
\hline
PILOT & Projeção numérica bidimensional; cinco tentativas de otimização \\
\hline
CLOISTER & Correlação 0,7; significância 0,05 \\
\hline
PYTHIA & SVM RBF; 30 configurações de $C$ e $\gamma$; quatro partições estratificadas; semente 0 \\
\hline
TRACE & \textit{Footprints} das classes previstas; pureza mínima 0,55 \\
\hline
\end{tabular}
\par\smallskip
{\fonte{elaboração própria a partir dos parâmetros adotados nos modelos.}}
\end{singlespace}
\end{quadro}

Nos modelos finais, o SIFTED filtrou pela magnitude da correlação de
Pearson significativa com desempenho, retendo também o atributo de maior
correlação absoluta por algoritmo. Os sobreviventes foram
agrupados por KMeans em dez grupos, e um algoritmo genético escolheu
um atributo por grupo, totalizando dez atributos. Cada combinação foi avaliada por PILOT e classificadores de
três vizinhos, em validação cruzada com cinco partições embaralhadas.
A aptidão foi o negativo do maior erro médio de classificação entre os
seis algoritmos; sua maximização buscou minimizar esse erro.

\newpage
Os parâmetros da seleção genética do Quadro~\ref{qua:met-matilda} foram
mantidos da configuração do SIFTED em pyInstanceSpace
\cite{PyInstanceSpace021} nos quatro ajustes, para aplicar o mesmo
procedimento de seleção de atributos. Não foram calibrados para o WSP nem
extraídos de \citeonline{SmithMilesBowly2015}, cuja busca genética gera
instâncias para pontos-alvo, em vez de selecionar atributos. Nesta
implementação, corrigiram-se localmente o filtro de correlação, a direção
da aptidão e o mapeamento dos atributos, sem alterar os valores do quadro.

A aptidão da combinação escolhida foi recalculada para conferir a seleção.
Nessa etapa, os indivíduos do algoritmo genético representaram combinações
de atributos; na geração da Seção~\ref{sec:met-geracao}, representaram
configurações de instâncias.

Após a projeção PILOT, PYTHIA padronizou as coordenadas e ajustou uma
SVM RBF por algoritmo. Uma busca aleatória de hiperparâmetros com
validação cruzada avaliou as configurações
de $C$ e $\gamma$ do Quadro~\ref{qua:met-matilda}, geradas por hipercubo
latino em $2^u$, $u\in[-10,4]$.
Os pesos corresponderam ao desvio absoluto
do desempenho em relação à média global dos \textit{gaps}; zeros foram substituídos
pelo menor peso positivo. O protocolo dos modelos finais utiliza quatro
partições estratificadas e embaralhadas, exigindo ao menos quatro exemplos
de cada classe por algoritmo em cada coleção. As cinco partições do mapa
de geração pertencem ao ajuste com rótulos históricos; os modelos finais
utilizam os rótulos do experimento comparativo atualizado.
CLOISTER estimou os limites, e TRACE construiu footprints das classes
previstas por PYTHIA, registrando área, densidade e pureza.

As predições finais na amostra foram auditadas contra os rótulos,
reconstruindo matrizes de confusão, acurácia, acurácia balanceada,
precisão e sensibilidade. Sem predições positivas, a precisão recebeu zero.
Como diagnóstico auxiliar interno não aninhado, realizaram-se dez repetições
de validação cruzada estratificada com quatro partições por modelo e
algoritmo, com sementes 0--9. Reutilizaram-se
$C$, $\gamma$ e coordenadas globais, sem a padronização adicional do
PYTHIA e com pesos de classes balanceados. Calcularam-se média e desvio de
acurácia balanceada, precisão, sensibilidade, F1, Matthews e área sob a curva
(\textit{area under the curve}, AUC) da característica de operação do receptor
(\textit{receiver operating characteristic}, ROC).
As repetições avaliaram a variação entre partições na representação ajustada.
A acurácia balanceada considerou as duas classes com igual peso,
reduzindo a influência da classe majoritária na interpretação.

Os contrastes controle/topologia e 55/74 compararam seleção pelo Jaccard,
ordenação das distâncias por Spearman, sobreposição dos cinco vizinhos e
deslocamento após Procrustes, com centralização, normalização e transformação
ortogonal com escala. Em 55/74, apenas os 55 identificadores comuns
participaram das distâncias e listas de vizinhos. Compararam-se também
predições binárias e footprints. A concordância mediu a coincidência das
predições entre modelos; a acurácia, os acertos em relação aos rótulos.
Os footprints foram interpretados na escala de cada mapa.
A sobreposição dos cinco vizinhos foi sensível a desempates numéricos e
não foi utilizada na interpretação dos resultados. Como diagnóstico
descritivo da associação com o tamanho, calcularam-se também correlações
de Spearman entre os quatro atributos topológicos selecionados no modelo
de 74 instâncias e as quantidades $|V|$ e $|A|$, usando os valores dos
atributos antes das transformações da ISA.

\section{Registro e análise das trajetórias do IBS}\label{sec:met-trajetorias}

Durante a execução do IBS, registraram-se
execução, objetivo e assinatura da solução devolvida ao fim de cada
iteração completa, adotada como unidade de amostragem da busca.
A construção seguiu o conceito de STN
\cite{ChaconSartoriEtAl2023,OchoaMalanBlum2021}, com dois protocolos:
\begin{alineas}
 \item histórico: 16 instâncias clássicas, dez sementes cada (160 execuções),
 $\beta=25$, $\eta=35$, $c=10$, até 30 segundos e dez iterações, com ótimo
 conhecido como alvo de parada; assinatura pelo conjunto ordenado dos
 vértices protegidos, sem instantes;
 \item adicional: 19 instâncias, dez sementes cada (190 execuções), inclusive
 Huge; parâmetros e limites do experimento principal, com assinatura de pares
 vértice--instante.
\end{alineas}

Preservaram-se os registros brutos. O filtro não piorante reteve, na ordem
original, registros com objetivo menor ou igual ao melhor já observado
na execução, incluindo empates. A rede histórica utilizou essa filtragem;
a análise adicional examinou tanto a sequência bruta quanto a filtrada.
A solução de uma iteração podia ser pior que a incumbente. Por isso,
guardaram-se tanto o último registro quanto a melhor solução acumulada.

Assinaturas iguais foram fundidas em nós. Transições entre assinaturas
diferentes formaram arestas direcionadas ponderadas pela frequência,
excluindo laços, isto é, transições de uma assinatura para ela mesma.
Essas repetições consecutivas não produziram arestas na rede;
as transições entre assinaturas distintas conservaram suas frequências. Calcularam-se nós, arestas e transições, assinaturas
de melhor objetivo observado e terminais fora desse conjunto. A recorrência
foi medida pela fração de nós compartilhados entre execuções e pelo Jaccard
médio entre pares de conjuntos visitados; registrou-se também a fração de
execuções terminadas no conjunto de melhor avaliação.

A força de entrada dos melhores nós somou os pesos recebidos e dividiu
pelo número de execuções. Nas redes clássicas, utilizou-se a razão limitada
a um; nas adições, calcularam-se a razão e sua versão limitada. A medida foi
interpretada como frequência normalizada de transições para os melhores nós.

Os descritores foram calculados em Python, e as figuras históricas foram
elaboradas em R. Nas adições, as assinaturas das soluções definiram os nós,
e os valores da função objetivo identificaram os nós de melhor avaliação
e os nós terminais fora desse conjunto. Para verificar a consistência das
trajetórias, compararam-se o melhor objetivo registrado, a solução final
e o histórico de melhorias de cada execução, mantendo sua identificação
por semente.

A integração utilizou o identificador da instância para associar os
descritores das STNs aos atributos e ao desempenho. Nas 16 clássicas,
cruzaram-se as redes filtradas com as coordenadas da PCA exploratória e
com o rótulo de melhor método da análise inicial, resumindo recorrência e
terminais por grupo de desempenho. Nas 19 adições, contrastaram-se os
perfis das redes com classe de tamanho e vantagem medida no experimento
comparativo. Essa análise descritiva reuniu os resultados das STNs, dos
atributos e do experimento. Os modelos MATILDA foram ajustados com os atributos
das instâncias, enquanto os descritores das trajetórias foram examinados
nessa comparação posterior.

\section{Tratamento e verificação dos resultados}\label{sec:met-analise}

A comparação utilizou soluções reavaliadas sob as mesmas restrições de
capacidade, disponibilidade de recursos e viabilidade temporal. Para obter
uma pontuação comum aos métodos, a propagação foi recalculada por Dijkstra
com os atrasos das alocações, contando os vértices atingidos estritamente
antes do horizonte. A análise utilizou essa pontuação e preservou os
registros originais. A associação por instância e semente permitiu
conferir a cobertura do plano experimental e evitar registros duplicados.

As duas execuções da ILS que falharam receberam o objetivo da solução sem
recursos e permaneceram na média principal, identificadas como falhas.
Uma análise de sensibilidade comparou \textit{gaps} e rótulos com e sem esses
registros. Na LBBD, distinguiram-se as soluções alternativas conforme o
status e a etapa executada, para interpretar os limites de acordo com a
formulação e a etapa que os produziram.

Seja $f_{i\alpha}$ a média dos objetivos validados nas $r_i$ sementes do
algoritmo $\alpha$ na instância $i$, e $b_i>0$ o menor objetivo validado
em qualquer execução dos seis métodos. O \textit{gap} percentual foi
\begin{equation}
 \mathit{gap}^{\mathrm{BKV}}_{i\alpha}
 =100\,\frac{f_{i\alpha}-b_i}{b_i}.
 \label{eq:gap-relativo}
\end{equation}
\noindent O BKV expressa a melhor observação validada no experimento.
A certificação de otimalidade exige adicionalmente um limite global válido.
Para uma coleção $C$, calculou-se
\begin{equation}
 G_{\alpha,C}=\frac{1}{|C|}\sum_{i\in C}
 \mathit{gap}^{\mathrm{BKV}}_{i\alpha}.
 \label{eq:met-agregado}
\end{equation}
\noindent Cada instância recebeu o mesmo peso, apesar das diferentes
quantidades de sementes. Médias por execução foram identificadas separadamente.

Por instância e algoritmo, extraíram-se média, mediana, extremos e desvio
padrão amostral, tempos e falhas. Os resumos abrangeram 55 originais,
19 adicionais, 74 totais e classes de tamanho. Contaram-se obtenções do
BKV, menores médias e vitórias exclusivas; empates de médias foram tratados
por igualdade dos valores calculados e receberam posições médias no ranking.
Origem, composição por classe e variabilidade acompanharam a interpretação.

A inferência utilizou as 74 médias pareadas por instância. Para os
15 pares de métodos, aplicou-se teste de sinal exato bicaudal, excluindo
empates do número efetivo de pares, e correção de Holm para a família
completa, com referência de 0,05. Contagens de vitórias, empates e derrotas
e diferenças de \textit{gap} acompanharam os testes, que avaliaram
a direção das diferenças na coleção dirigida.
\textit{Gaps} internos dos \textit{solvers} foram analisados separadamente do \textit{gap} empírico.

A integração manteve a associação entre cada instância, seus atributos,
os desempenhos e as trajetórias. As tabelas e figuras foram calculadas a
partir dos resultados correspondentes, preservando a identificação do
experimento, do modelo e do protocolo de trajetórias. Essa associação permitiu
interpretar conjuntamente as etapas sem confundir seus diferentes protocolos.

\section{Escopo e limites da análise}\label{sec:met-limitacoes}

O povoamento foi avaliado pela cobertura do domínio interno do mapa
fixado para a geração. A medida descreve a ocupação dessa representação;
a qualidade preditiva foi examinada pelos diagnósticos dos modelos finais.
A geração não incluiu controle aleatório sob igual orçamento;
o aumento de cobertura não demonstra superioridade sobre geração aleatória.
A divisão do orçamento entre evolução e busca local não foi otimizada para o WSP.
A validade das adições foi estabelecida pela consistência com o gerador e com o modelo WSP.

A comparação entre métodos adotou parâmetros comuns às classes e
orçamentos proporcionais ao tamanho das grades. Origem e tamanho foram
considerados conjuntamente, pois 15 das 19 adições são Large ou Huge.
A análise procura associações entre atributos e desempenho; a identificação
de efeitos individuais exigiria variações controladas dessas propriedades.
O orçamento por tempo avalia o conjunto formado pela estratégia de busca,
pela implementação e por sua configuração. As diferenças de linguagem
também participam desse resultado.
A precisão das médias em Huge depende das três sementes utilizadas;
a distribuição complementar dos trabalhos entre máquinas não constituiu
um teste do efeito do equipamento.

Os modelos ISA têm limites de interpretação. O contraste entre atributos
originais e topológicos utilizou um único critério de bom desempenho:
\textit{gap} de até 10\%. As classes e regiões dependem desse limiar, e
os experimentos registrados não avaliaram sua sensibilidade a outros
valores. Além disso, há muitos candidatos, 107, para 55 ou 74 instâncias,
e classes desbalanceadas, o que pode tornar os classificadores sensíveis
à partição dos dados. Os diagnósticos descrevem essa configuração e sua
variação entre partições.

A seleção foi executada com uma única semente. As repetições de validação
cruzada não medem a estabilidade da seleção sob outras sementes do SIFTED.
Assim, as diferenças entre controle e combinação topológica descrevem os
ajustes obtidos nesta configuração, sem estabelecer um efeito estável
ou universal da topologia.

Transformações, seleção, projeção e hiperparâmetros foram estimados com
todas as instâncias da respectiva coleção, sem repetir o fluxo completo
separadamente nos dados de ajuste de cada partição. As métricas são,
portanto, diagnósticos internos e podem ser otimistas; não demonstram
generalização para instâncias novas. Essa avaliação exigiria dados externos
ou validação aninhada, com todas as etapas ajustadas dentro das partições.

As STNs históricas e adicionais foram interpretadas separadamente porque
seus protocolos diferem na assinatura da solução, na parada e nos parâmetros
(Seção~\ref{sec:met-trajetorias}). As redes não são diretamente comparáveis,
e suas associações com atributos e desempenho têm caráter descritivo.
